% STATE: DRAFT (Section 5 awaiting human review).
\section{Examples and Evaluation}\label{sec:evaluation}

\subsection{Worked Examples}

% PROV[user]: Approved request/grant, two-thread lock and traffic-light examples
% with reserved illustration space and TikZ artist briefs.
% PROV[code]: Examples/Demos/Walkthrough.lean; Examples/Demos/Frontend.lean;
% Examples/Models/TwoThreadLock.lean; Examples/Models/TrafficLight.lean.
The examples exercise different parts of the interface and proof pipeline.
Request/grant follows a formula through checking and executable extraction.
The lock model contrasts safety with starvation, while the traffic light
compares direct semantic proofs with supplied-monitor checks. All use the
infinite-path semantics described in \cref{sec:correctness}.

\subsubsection{Request and Grant: a Complete User Session}

% PROV[code]: Examples/Demos/Walkthrough.lean:L30-L78,L114-L186;
% Examples/Demos/Frontend.lean:L36-L50.
The system in \sourcefile{Examples/Demos/Walkthrough.lean} has the two
variants shown in \cref{fig:example-request-grant}. Both are serial, but the
waiting self-loop permits an infinite execution without a grant. Only the
healthy variant satisfies $\mathbf{F}\,\mathit{grant}$ on every initial
path. The concrete true and false verdicts are proved with
\lean{kilt_decide}; \sourcefile{Examples/Demos/Frontend.lean} additionally
proves healthy satisfaction with \lean{kilt_verify}.

% PROV[code]: Examples/Demos/Walkthrough.lean:L157-L186,L192-L235.
% PROV[build]: 2026-11-15, ./scripts/verify-examples.sh --verbose, exit 0;
% .lake/logs/build-20261115T140749Z-0o5nGo.log, walkthrough output.
The displayed healthy result has no counterexample. The waiting variant
returns a system lasso with empty prefix and loop $[\mathit{waiting}]$:
granting remains possible but is never chosen. The theorem
\lean{returned_path_violates} concerns the returned lasso's own infinite
projection. Executed assertions also check the same verdict/witness outcomes
for $\mathbf{G}(\mathit{request}\rightarrow\mathbf{F}\,\mathit{grant})$.
These response displays are execution checks; the named verdict and
returned-path theorems above concern the eventual-grant formula.

% PROV[user]: Supplied report/assets/example-request-grant.tex for integration.
% PROV[code]: Examples/Demos/Walkthrough.lean, model, label and display.
\begin{figure}[!ht]
  \centering
  \resizebox{\linewidth}{!}{% example-request-grant.tex  --  Section 5.1, request/grant example.
% Paired-figure fragment (about 158 mm x 44 mm): no preamble, caption or figure
% environment. Needs only TikZ and the kilt* styles/colours of report/src/preamble.tex.
% Facts follow Examples/Demos/Walkthrough.lean (State, label, model, eventuallyGrant).
%   (a) two states, two variants; only "may-wait" has waiting -> waiting
%   (b) the displayed system lasso for F grant: prefix = [], loop = [waiting]
\begin{tikzpicture}[x=1mm, y=1mm,
  rgstate/.style={kiltstage, circle, draw=kiltnavy, fill=kiltpale, line width=0.8pt,
    minimum size=14mm, inner sep=0pt, font=\footnotesize\sffamily, text=kiltnavy},
  rgviol/.style={rgstate, draw=kiltorange, fill=kiltpeach, text=kiltorange,
    line width=1.2pt, minimum size=13mm},
  rgprop/.style={font=\footnotesize, text=kiltblue, inner sep=1pt},
  rgtext/.style={font=\footnotesize, text=kiltnavy, inner sep=0pt},
  rgtitle/.style={font=\footnotesize\bfseries, text=kiltnavy, anchor=west, inner sep=0pt},
  rgcex/.style={kiltarrow, draw=kiltorange, line width=1.3pt},
]

% =====================================================================
% (a)  one model, two variants                         x = 0 .. 76 mm
% =====================================================================
\node[rgtitle] at (0,43) {(a) Two variants of one model};

\node[rgstate] (w) at (22,21) {waiting};
\node[rgstate] (g) at (56,21) {granted};

\draw[kiltarrow] (3,21) -- (w.west);                                  % initial
\draw[kiltarrow] (w) -- (g);                                          % waiting -> granted
\draw[kiltarrow] (g.125) to[out=125,in=55,looseness=4] (g.55);      % granted -> granted
\draw[kiltarrow, kiltdashed, kiltorange, line width=1.2pt]
  (w.125) to[out=125,in=55,looseness=4] (w.55);                     % waiting -> waiting (may-wait only)
\node[font=\footnotesize, text=kiltorange, inner sep=0pt] at (22,37.2) {may-wait only};

\node[rgprop] at (22,11.2) {$\{\mathit{request}\}$};
\node[rgprop] at (56,11.2) {$\{\mathit{grant}\}$};

% legend + verdicts for F grant
\draw[kiltnavy, line width=0.8pt] (0.5,6.3) -- (7.5,6.3);
\node[rgtext, anchor=west] at (9.5,6.3) {healthy: solid edges only};
\node[rgtext, anchor=east] at (76,6.3) {$\mathbf{F}\,\mathit{grant}$: true};
\draw[kiltdashed, kiltorange, line width=1.2pt] (0.5,2.3) -- (7.5,2.3);
\node[rgtext, anchor=west] at (9.5,2.3) {may-wait: adds dashed loop};
\node[rgtext, anchor=east, text=kiltorange] at (76,2.3) {$\mathbf{F}\,\mathit{grant}$: false};

% =====================================================================
% (b)  displayed lasso for F grant (may-wait variant)   x = 82 .. 158 mm
% =====================================================================
\begin{scope}[shift={(82,0)}]
\node[rgtitle] at (0,43) {(b) Displayed lasso for $\mathbf{F}\,\mathit{grant}$ (may-wait)};

\node[rgviol] (e1) at (10.5,27) {waiting};
\node[rgviol] (e2) at (29.5,27) {waiting};
\node[rgviol] (e3) at (48.5,27) {waiting};
\draw[rgcex] (e1) -- (e2);
\draw[rgcex] (e2) -- (e3);
% the execution never terminates
\draw[rgcex, dotted, line cap=round, dash pattern=on 0pt off 2.5pt] (e3.east) -- (74,27);

\node[rgprop] at (10.5,35.2) {$\{\mathit{request}\}$};
\node[rgprop] at (29.5,35.2) {$\{\mathit{request}\}$};
\node[rgprop] at (48.5,35.2) {$\{\mathit{request}\}$};

% brace under the repeated states: "repeat forever"
\draw[kiltorange, line width=0.8pt]
  (4,19) to[out=-90,in=180] (5.5,17.5) -- (38,17.5) to[out=0,in=90] (39.5,16)
         to[out=90,in=180] (41,17.5)   -- (73.5,17.5) to[out=0,in=-90] (75,19);
\node[font=\footnotesize, text=kiltorange, inner sep=0pt] at (39.5,12.6) {repeat forever};

\node[rgtext, anchor=west] at (0,7.4) {$\pi=\textsf{waiting}^{\omega}$};
\node[rgtext, anchor=east, text=kiltorange] at (76,7.4) {$\mathit{grant}$ never holds};
\node[rgtext, anchor=west, text=kiltblue] at (0,2.4)
  {$\mathit{grant}$ remains possible: \textsf{waiting}$\,\to\,$\textsf{granted}, unused};
\end{scope}

\end{tikzpicture}
}
  \caption{Request/grant: one added self-loop permits an infinite violation.}
  \label{fig:example-request-grant}
\end{figure}

\subsubsection{Two Threads: Mutual Exclusion and Starvation}

% PROV[code]: Examples/Models/TwoThreadLock.lean:L3-L16,L28-L90,L114-L168.
The model in \sourcefile{Examples/Models/TwoThreadLock.lean} has two thread
program counters and one shared lock bit. Each transition schedules one
logical instruction of either thread; reads and writes observe that shared
bit. In the broken version, testing and claiming the lock are separate
steps. Both threads can therefore pass the test before either claims it.
\lean{race_collision} establishes that both are critical at position six of
the explicit execution, refuting
$\mathbf{G}\neg(\mathit{critical}_0\land\mathit{critical}_1)$.
The atomic variant combines test and claim; an invariant proves mutual
exclusion along every initial execution.

% PROV[code]: Examples/Models/TwoThreadLock.lean:L179-L233.
Without fairness, atomicity still permits starvation. An explicit schedule
lets thread zero request once, then schedules only thread one forever.
Thread zero remains waiting, violating
$\mathbf{G}(\mathit{request}_0\rightarrow
\mathbf{F}\,\mathit{critical}_0)$.
The false supplied-monitor verdict is derived from this semantic
counterexample and checker correctness; a separate \texttt{\#eval}
executes the search.

% PROV[user]: Supplied report/assets/example-two-thread-lock.tex for integration.
% PROV[code]: Examples/Models/TwoThreadLock.lean, raceSchedule and unfairSchedule.
\begin{figure}[!ht]
  \centering
  \resizebox{\linewidth}{!}{% example-two-thread-lock.tex  --  Section 5.1, two-thread lock example.
% Follows Examples/Models/TwoThreadLock.lean (raceSchedule, race_collision,
% unfairSchedule, unfair_waits). These are explicit semantic constructions,
% not output of the formula compiler or lasso extractor.
%   state = (pc0, pc1, L);  I idle, T trying, C checked, K critical.
%   (a) finite race: both threads test while L = 0, then both claim.
%   (b) atomic test + claim: mutual exclusion holds, thread 0 starves.

\begin{tikzpicture}[x=1mm, y=1mm,
  ltbox/.style={kiltstage, draw=kiltnavy, fill=kiltpale, line width=0.8pt,
    minimum width=17mm, minimum height=4.2mm, inner sep=0pt,
    font=\footnotesize\sffamily, text=kiltnavy},
  lttest/.style={ltbox, draw=kiltblue, line width=1.6pt},
  ltbad/.style={ltbox, draw=kiltorange, fill=kiltpeach, text=kiltorange, line width=1.5pt},
  ltloop/.style={ltbox, minimum width=14mm, minimum height=5mm,
    draw=kiltorange, fill=kiltpeach, text=kiltorange, line width=1.2pt},
  ltpos/.style={font=\footnotesize, text=kiltblue, inner sep=0pt},
  lttxt/.style={font=\footnotesize, text=kiltnavy, inner sep=0pt, anchor=west},
  ltnote/.style={font=\footnotesize, inner sep=0pt, anchor=west},
  lttitle/.style={font=\footnotesize\bfseries, text=kiltnavy, anchor=west, inner sep=0pt},
  ltcex/.style={kiltarrow, draw=kiltorange, line width=1.3pt},
]

% =====================================================================
% (a)  both test before either claims                  x = 0 .. 76 mm
% =====================================================================
\node[lttitle] at (0,56) {(a) Test, then claim: no retest};
\node[ltnote, text=kiltorange] at (0,52.4) {mutual exclusion fails at position 6};

% position | state | transition into this position
\node[ltpos] at (2.6,47)   {0}; \node[ltbox]  (p0) at (14,47)   {(I,\,I,\,0)};
\node[lttxt] at (26,47)   {initial};
\node[ltpos] at (2.6,40.8) {1}; \node[ltbox]  (p1) at (14,40.8) {(T,\,I,\,0)};
\node[lttxt] at (26,40.8) {$T_0$ requests};
\node[ltpos] at (2.6,34.6) {2}; \node[lttest] (p2) at (14,34.6) {(C,\,I,\,0)};
\node[lttxt, text=kiltblue, font=\footnotesize\bfseries] at (26,34.6) {$T_0$ tests the free lock};
\node[ltpos] at (2.6,28.4) {3}; \node[ltbox]  (p3) at (14,28.4) {(C,\,T,\,0)};
\node[lttxt] at (26,28.4) {$T_1$ requests};
\node[ltpos] at (2.6,22.2) {4}; \node[lttest] (p4) at (14,22.2) {(C,\,C,\,0)};
\node[lttxt, text=kiltblue, font=\footnotesize\bfseries] at (26,22.2) {$T_1$ tests the free lock};
\node[ltpos] at (2.6,16)   {5}; \node[ltbox]  (p5) at (14,16)   {(K,\,C,\,1)};
\node[lttxt] at (26,16)   {$T_0$ claims};
\node[ltpos, text=kiltorange] at (2.6,9.8) {6}; \node[ltbad] (p6) at (14,9.8) {(K,\,K,\,1)};
\node[lttxt, text=kiltorange, font=\footnotesize\bfseries] at (26,9.8) {$T_1$ claims (no retest)};

\foreach \a/\b in {p0/p1, p1/p2, p2/p3, p3/p4, p4/p5, p5/p6}
  \draw[kiltarrow] (\a.south) -- (\b.north);
% the execution goes on after this finite safety witness
\draw[kiltarrow, kiltdashed, kiltblue] (p6.south) -- (14,4.2);

% =====================================================================
% (b)  atomic acquisition still permits starvation      x = 82 .. 158 mm
% =====================================================================
\begin{scope}[shift={(82,0)}]
\node[lttitle] at (0,56) {(b) Atomic lock, unfair schedule};
\node[ltnote, text=kiltblue] at (0,52.4) {test + claim in one step};

\node[ltbox, minimum width=14mm, minimum height=5mm] (i0) at (13,40) {(I,\,I,\,0)};
\node[ltloop] (a) at (38,40)   {(T,\,I,\,0)};
\node[ltloop] (b) at (65,40)   {(T,\,T,\,0)};
\node[ltloop] (c) at (51.5,18.5) {(T,\,K,\,1)};

\draw[kiltarrow] (0.5,40) -- (i0.west);                         % initial
\draw[kiltarrow] (i0) -- (a);                                   % thread 0, once
\draw[ltcex] (a) -- (b);                                        % thread 1 requests
\draw[ltcex] (b) -- (c);                                        % thread 1 test + claim
\draw[ltcex] (c) -- (a);                                        % thread 1 releases

\node[ltnote, anchor=south, text=kiltblue] at (25.5,44) {$T_0$};
\node[ltnote, anchor=north, text=kiltnavy] at (25.5,37.8) {prefix};
\node[ltnote, anchor=south, text=kiltorange] at (51.5,41.8) {request};
\node[ltnote, anchor=west, text=kiltorange] at (60.5,29) {acquire};
\node[ltnote, anchor=east, text=kiltorange] at (40.5,29) {release};
\node[ltnote, anchor=center, text=kiltorange, align=center] at (51.5,31) {repeats\\forever};

\node[ltnote, text=kiltblue]   at (0,12.2) {mutual exclusion holds};
\node[ltnote, text=kiltorange] at (0,8.3)  {$T_0$ can starve: no fairness};
\node[ltnote, text=kiltnavy]   at (0,4.4)  {all loop edges schedule $T_1$};
\end{scope}

% shared legend
\node[font=\footnotesize, text=kiltnavy, inner sep=0pt] at (79,0.2)
  {state $=(\mathit{pc}_0,\mathit{pc}_1,L)$, $L$ the lock bit;\quad
   I idle, T trying, C checked, K critical};

\end{tikzpicture}
}
  \caption{The lock example separates a finite safety violation from an
    infinite starvation execution.}
  \label{fig:example-lock}
\end{figure}

\subsubsection{Traffic Light: Safety Does Not Force Progress}

% PROV[code]: Examples/Models/TrafficLight.lean:L28-L85,L144-L173,L191-L244.
The controller in \sourcefile{Examples/Models/TrafficLight.lean} starts at
\lean{stop} and cycles through \lean{go} and \lean{caution}, with labels red,
green and yellow respectively. Its second variant adds a \lean{stop}
self-loop. Direct semantic proofs show that exactly one colour is lit on
every path of either model, and that the normal controller respects the
next-colour order. The extra loop preserves colour safety but permits red
forever, violating both $\mathbf{F}\,\mathit{green}$ and
$\mathbf{G}(\mathit{red}\rightarrow\mathbf{F}\,\mathit{green})$.
The four supplied-monitor verdicts (two properties on two models) are proved
by kernel computation. Language theorems transfer them to the declared
semantics; the permanent-red path is also constructed explicitly.
These calls use \lean{checkAgainst}, rather than automatic formula compilation.

% PROV[user]: Supplied report/assets/example-traffic-light.tex for integration.
% PROV[code]: Examples/Models/TrafficLight.lean, controller, stuckController and stuckPath.
\begin{figure}[!ht]
  \centering
  \resizebox{\linewidth}{!}{% example-traffic-light.tex  --  Section 5.1, traffic-light example.
% Follows Examples/Models/TrafficLight.lean (label, controller, stuckController,
% safety, eventuallyGreen, redResponds, stuckPath). Nodes are controller states;
% {red}, {green}, {yellow} are proposition labels and not automaton states.
%   (a) normal cycle   
%   (b) the same cycle plus the single extra edge stop -> stop

\begin{tikzpicture}[x=1mm, y=1mm,
  tlstate/.style={kiltstage, circle, draw=kiltnavy, fill=kiltpale, line width=0.8pt,
    minimum size=13mm, inner sep=0pt, font=\footnotesize\sffamily, text=kiltnavy},
  tlprop/.style={font=\footnotesize, text=kiltblue, inner sep=1pt, anchor=west},
  tlblock/.style={font=\footnotesize, text=kiltnavy, inner sep=0pt, align=left,
    anchor=north west, text width=39mm},
  tltitle/.style={font=\footnotesize\bfseries, text=kiltnavy, anchor=west, inner sep=0pt},
]

% =====================================================================
% (a)  normal labelled cycle                           x = 0 .. 76 mm
% =====================================================================
\node[tltitle] at (0,41) {(a) Normal labelled cycle};

\node[tlstate] (s) at (12.5,18)  {stop};
\node[tlstate] (g) at (30.5,30.5) {go};
\node[tlstate] (c) at (30.5,6.5)  {caution};

\draw[kiltarrow] (0.3,18) -- (s.west);          % initial
\draw[kiltarrow] (s) -- (g);
\draw[kiltarrow] (g) -- (c);
\draw[kiltarrow] (c) -- (s);

\node[tlprop, anchor=north] at (12.5,10.5) {$\{\mathit{red}\}$};
\node[tlprop] at (37.5,30.5) {$\{\mathit{green}\}$};
\node[tlprop] at (37.5,6.5)  {$\{\mathit{yellow}\}$};

\node[tlblock] at (38,26.5)
  {exactly one colour\\
   eventual green holds\\
   red response holds};

% =====================================================================
% (b)  the added self-loop                              x = 82 .. 158 mm
% =====================================================================
\begin{scope}[shift={(82,0)}]
\node[tltitle] at (0,41) {(b) One added self-loop};

\node[tlstate, draw=kiltorange, fill=kiltpeach, text=kiltorange] (s) at (12.5,18) {stop};
\node[tlstate] (g) at (30.5,30.5) {go};
\node[tlstate] (c) at (30.5,6.5)  {caution};

\draw[kiltarrow] (0.3,18) -- (s.west);          % initial
\draw[kiltarrow, kiltblue!40] (s) -- (g);        % unchosen edge: faded, still visible
\draw[kiltarrow] (g) -- (c);
\draw[kiltarrow] (c) -- (s);
\draw[kiltarrow, kiltdashed, kiltorange, line width=1.3pt]
  (s.150) to[out=150,in=85,looseness=4.5] (s.85);                      % stop -> stop

\node[font=\footnotesize, text=kiltorange, inner sep=0pt] at (7,36) {added loop};

\node[tlprop, anchor=north] at (12.5,10.5) {$\{\mathit{red}\}$};
\node[tlprop] at (37.5,30.5) {$\{\mathit{green}\}$};
\node[tlprop] at (37.5,6.5)  {$\{\mathit{yellow}\}$};

\node[tlblock] at (38,26.5)
  {$\pi=\textsf{stop}^{\omega}$\\
   one-colour safety holds\\
   \textcolor{kiltorange}{eventual green fails}\\
   \textcolor{kiltorange}{red response fails}};
\end{scope}

\node[font=\footnotesize, text=kiltblue, inner sep=0pt] at (79,-4)
  {Both variants are serial; no fairness is assumed.};

\end{tikzpicture}
}
  \caption{A serial controller can remain safe while failing liveness.}
  \label{fig:example-traffic}
\end{figure}

\clearpage
\subsection{Verification and Deliverable Status}

% PROV[code]: Examples.lean; Examples/Regression/BuchiTrust.lean:L27-L37;
% scripts/verify-examples.sh; lean-toolchain.
% PROV[build]: 2026-11-15, LEAN_NUM_THREADS=1 ./scripts/verify-examples.sh --verbose,
% exit 0, 1857 jobs, 405 axiom-audit messages, no warning/error diagnostics;
% .lake/logs/build-20261115T140749Z-0o5nGo.log.
Validation used the pinned Lean toolchain with the command
\begin{quote}\small
\begin{verbatim}
LEAN_NUM_THREADS=1 ./scripts/verify-examples.sh --verbose
\end{verbatim}
\end{quote}
The examples target and its library dependencies built successfully on
15 November 2026, with no warnings and 405 passing transitive axiom audits.
This was an incremental build: Lake can reuse up-to-date elaboration and
execution results. It is build evidence, not a performance measurement.

% PROV[code]: Examples/Regression/BuchiTrust.lean:L27-L37 and declaration list;
% KILT/LTL/Compile.lean; KILT/Kripke/Check.lean; KILT/Kripke/Counterexample.lean.
% PROV[build]: 2026-11-15, rg for sorry/custom axiom declarations in KILT and Examples
% returned no matches; direct #print axioms checks on 15 selected example results
% each reported exactly [propext, Classical.choice, Quot.sound], exit 0.
The audit permits only \lean{propext}, \lean{Classical.choice} and
\lean{Quot.sound} in its named results and fails on unfinished-proof
dependencies or additional axioms. Direct inspections of selected example
theorems reported exactly these three axioms. Source scans found no \lean{sorry} or
custom axiom declarations in \lean{KILT/} or \lean{Examples/}.
In \cref{tab:deliverable-status}, \emph{verified} refers to the cited Lean
contracts under their stated assumptions. Concrete execution checks
exercise examples and regressions; they do not replace those contracts.

% PROV[user]: Proposal-aligned deliverables approved in Section 1 and requested
% verification/deliverable table for Section 5.2.
\begin{table}[!ht]
  \centering
  \small
  \renewcommand{\arraystretch}{1.22}
  \begin{tabularx}{\linewidth}{@{}
    >{\raggedright\arraybackslash}p{0.19\linewidth}
    >{\raggedright\arraybackslash}X
    >{\raggedright\arraybackslash}p{0.25\linewidth}@{}}
    \toprule
    Deliverable & Principal source & Status in this validation \\
    \midrule
    % PROV[code]: KILT/LTL/Syntax.lean; KILT/LTL/Semantics.lean; KILT/LTL/Basic.lean.
    LTL basis &
      \sourcefile{LTL/Syntax.lean}, \sourcefile{LTL/Semantics.lean};
      laws in \sourcefile{LTL/Basic.lean} &
      Defined; semantic laws checked \\
    % PROV[code]: KILT/LTL/TableauCorrect.lean; KILT/LTL/Compile.lean;
    % KILT/Automata/GeneralizedBuchiCorrect.lean; current audit log.
    Formula compilation &
      \sourcefile{LTL/TableauCorrect.lean}, \sourcefile{LTL/Compile.lean},
      \sourcefile{Automata/GeneralizedBuchiCorrect.lean} &
      Verified language equivalences \\
    % PROV[code]: KILT/Automata/Product.lean; KILT/Graph/Acceptance.lean;
    % KILT/Graph/Buchi.lean; current audit log.
    Product and emptiness &
      \sourcefile{Automata/Product.lean}, \sourcefile{Graph/Acceptance.lean},
      \sourcefile{Graph/Buchi.lean} &
      Verified run/cycle and search correspondences \\
    % PROV[code]: KILT/Kripke/Check.lean; KILT/Kripke/PreparedCheck.lean; current audit log.
    Public checker &
      \sourcefile{Kripke/Check.lean},
      \sourcefile{Kripke/PreparedCheck.lean} &
      Verified verdict equivalence \\
    % PROV[code]: KILT/Automata/ProductLasso.lean; KILT/Kripke/Counterexample.lean;
    % current audit log.
    Lasso extraction &
      \sourcefile{Automata/ProductLasso.lean},
      \sourcefile{Kripke/Counterexample.lean} &
      Verified returned trace and witness presence \\
    % PROV[code]: KILT/Frontend.lean; Examples/Demos; Examples/Models;
    % Examples/Regression/FrontendRegression.lean; current examples build.
    User interface and examples &
      \sourcefile{Frontend.lean}; \lean{Examples/Demos/},
      \lean{Examples/Models/}, \lean{Examples/Regression/} &
      Proof macros and concrete checks pass \\
    \bottomrule
  \end{tabularx}
  \caption{Implemented deliverables and their checked contracts.
    Source paths are relative to \lean{KILT/} unless an \lean{Examples/} path is shown.}
  \label{tab:deliverable-status}
\end{table}

% PROV[user]: Section-by-section report review and deferred final submission.
% PROV[code]: scripts/verify-examples.sh and scripts/compile-report.sh are separate targets.
The report and illustrations remain under review. This validation covers
Lean targets; HTML documentation generation and final report approval are
separate deliverables.

\clearpage
\subsection{Experimental Results and Complexity}\label{sec:experimental-cost}

% PROV[user]: Approved the archived experiment results and requested a rough
% complexity analysis, keeping measured results distinct from estimates.
% Requested plain wording focused on the implementation, with benchmark
% limitations stated together rather than repeated throughout.
% Requested describing a one-off test repeatable against older commit versions,
% without a visible citation to the benchmark archive.
% PROV[code]: .ai/archive/benchmarks-2026-10-27-p3-p5/RESULTS.md;
% baseline/environment.json; p5/environment.json; p5/source/Benchmarks/GraphStorage.lean.
The one-off test in \cref{tab:historical-performance} compared versions on
27 October 2026. The optimised version uses an array-based
discovery queue, marks accepting states over the reachable graph, and checks
source literals before enumerating tableau successors. A common script ran
each version three times as a native executable on an x86-64 Linux virtual
machine, pinned to one virtual CPU with one Lean thread. Executable
compilation time was excluded. The timed calls include LTL translation and
product preparation. The comparison can be repeated by checking out older
commit versions and running the same test script.

% PROV[code]: baseline/metrics.json and p5/metrics.json in the experiment archive.
% PROV[build]: 2026-11-16, extracted four native sample arrays per snapshot with jq;
% recomputed each median by numerical sorting and taking the middle of three;
% rounded to two decimal places. No new benchmark execution.
\begin{table}[!ht]
  \centering
  \small
  \renewcommand{\arraystretch}{1.15}
  \begin{tabular}{@{}lrrrr@{}}
    \toprule
    & \multicolumn{2}{c}{Check (ms)}
    & \multicolumn{2}{c}{Extract counterexample (ms)} \\
    Workload & Baseline & Optimised & Baseline & Optimised \\
    \midrule
    Waiting-response & 21.54 & 10.19 & 21.84 & 10.29 \\
    Traffic-response & 25.13 & 10.99 & 24.67 & 11.66 \\
    \bottomrule
  \end{tabular}
  \caption{Median execution times before and after optimisation, three runs per version.}
  \label{tab:historical-performance}
\end{table}

% PROV[code]: Experiment RESULTS.md, Variability and follow-up;
% paired-native/RESULTS.md; KILT/LTL/Tableau.lean:L159-L178.
The main checks took about half as long after the changes. The successor
function skips candidate enumeration when the label is incompatible, and
the prepared graph stores each discovered state's successors for reuse. The versions returned
the same verdicts and selected witnesses. A further comparison immediately
before and after the literal-check change used eight pairs of native runs;
both main checks were faster in every pair. The compatible-label test was
$6.36\%$ slower by median paired change, while sparse-product preparation was
$9.22\%$ slower. The cause of the latter increase was not isolated.

% PROV[code]: Experiment RESULTS.md, methodology and Variability and follow-up;
% paired-native/RESULTS.md, limitations of the short-duration comparisons.
These measurements describe the tested models on this virtual machine.
The baseline briefly overlapped separate proof checks, and the short-duration
tests were sensitive to variation between runs.

% PROV[code]: KILT/LTL/Compile.lean:L89-L98, card_automatonState_le;
% KILT/LTL/Tableau.lean, State and successors; KILT/Automata/Product.lean:L21-L29.
% PROV[build]: card_automatonState_le passed the standard-axiom audit in
% .lake/logs/build-20261115T140749Z-0o5nGo.log.
For the size bound, let $Q$ be the compiled automaton's state set, $n$ the
number of subformulas after normalising the negated property, $m$ its until
obligations and $N$ the number of model states.
With $k=\max(1,m)$, \sourcefile{LTL/Compile.lean} proves the automaton bound below; the product
bound follows by counting pairs:
\[
  |Q|\le k\,2^n,\qquad V\le N|Q|\le Nk\,2^n.
\]
Here $V$ counts reachable product vertices. The compiler enumerates subsets
of the formula's subformulas, giving exponential growth in the worst case.
Literal filtering saves enumeration work when the input label is incompatible.

% PROV[code]: KILT/Graph/Basic.lean, FiniteGraph.row and edge;
% KILT/Graph/Traversal.lean, Visits.advance and Visits.search;
% KILT/Graph/Buchi.lean, PreparedReachable.hasAcceptingCycle;
% KILT/Graph/Internal/Parents.lean, Parents and unwind;
% KILT/Graph/Prepare.lean, Buffer.advance and prepareFromWork.
% PROV[code-derived estimate]: Frontiers contain newly visited vertices only.
% Across a return search, each source enters a frontier once; checking its
% adjacency row for at most V targets costs O(V*E) in total. Per-level array
% scans contribute O(V^2). Assumes unit-cost dense-ID equality/array access;
% excludes compiler/preparation costs and is not a Lean-proved runtime theorem.
For a rough operation count, assume constant-cost comparisons of vertex IDs
and array access. A prepared graph with $E$ edges stores $O(V+E)$ graph
entries; visited or parent arrays add $O(V)$. Frontier search in
\sourcefile{Graph/Traversal.lean} scans possible target vertices and tests
membership in stored successor rows. Each vertex enters a frontier once,
and its row can be tested against up to $V$ targets, giving a return-query
estimate of $O(V(V+E))$. With $a\le V$ reachable accepting vertices, the
checker can perform up to $a$ such queries. Preparation additionally includes
formula-state enumeration, state-order lookups and finite-set operations.
\clearpage
